% OpenStax Calculus Volume 2 -- Section 2.4 Exercises: Arc Length of a Curve and Surface Area
% Exercises reproduced from OpenStax, licensed under CC BY 4.0.
% Source: https://openstax.org/books/calculus-volume-2/pages/2-4-arc-length-of-a-curve-and-surface-area
\documentclass{exercises}
\title{OpenStax Calculus Volume 2}
\subtitle{Section 2.4 Exercises: Arc Length of a Curve and Surface Area}
\sourceurl{https://openstax.org/books/calculus-volume-2/pages/2-4-arc-length-of-a-curve-and-surface-area}
\author{}
\date{}

\begin{document}
\maketitle


For the following exercises, find the length of the functions over the given interval.

\begin{enumerate}[start=1]
\item $y=5x$ from $x=0$ to $x=2$
\item $y=-\frac{1}{2}x+25$ from $x=1$ to $x=4$
\item $x=4y$ from $y=-1$ to $y=1$
\item Pick an arbitrary linear function $x=g(y)$ over any interval of your choice $(y_{1},y_{2})$. Determine the length of the function and then prove the length is correct by using geometry.
\item Find the surface area of the volume generated when the curve $y=\sqrt{x}$ revolves around the $x\text{-axis}$ from $(1,1)$ to $(4,2)$, as seen here. \\ \textit{[Figure: This figure is a surface. It has been formed by rotating the curve y=squareroot(x) about the x-axis. The surface is inside of a cube to show 3-dimensions.]}
\item Find the surface area of the volume generated when the curve $y=x^{2}$ revolves around the $y\text{-axis}$ from $(1,1)$ to $(3,9)$. \\ \textit{[Figure: This figure is a surface. It has an elliptical shape to the top, forming a ``bowl''.]}
\end{enumerate}

For the following exercises, find the lengths of the functions of $x$ over the given interval. If you cannot evaluate the integral exactly, use technology to approximate it.

\begin{enumerate}[resume]
\item $y=x^{3/2}$ from $(0,0)$ to $(1,1)$
\item $y=x^{2/3}$ from $(1,1)$ to $(8,4)$
\item $y=\frac{1}{3}(x^{2}+2)^{3/2}$ from $x=0$ to $x=1$
\item $y=\frac{1}{3}(x^{2}-2)^{3/2}$ from $x=2$ to $x=4$
\item \techrequired $y=e^{x}$ on $x=0$ to $x=1$
\item $y=\frac{x^{3}}{3}+\frac{1}{4x}$ from $x=1$ to $x=3$
\item $y=\frac{x^{4}}{4}+\frac{1}{8x^{2}}$ from $x=1$ to $x=2$
\item $y=\frac{2x^{3/2}}{3}-\frac{x^{1/2}}{2}$ from $x=1$ to $x=4$
\item $y=\frac{1}{27}(9x^{2}+6)^{3/2}$ from $x=0$ to $x=2$
\item \techrequired $y=\sin x$ on $x=0$ to $x=\pi$
\end{enumerate}

For the following exercises, find the lengths of the functions of $y$ over the given interval. If you cannot evaluate the integral exactly, use technology to approximate it.

\begin{enumerate}[resume]
\item $y=\frac{5-3x}{4}$ from $y=0$ to $y=4$
\item $x=\frac{1}{2}(e^{y}+e^{-y})$ from $y=-1$ to $y=1$
\item $x=5y^{3/2}$ from $y=0$ to $y=1$
\item \techrequired $x=y^{2}$ from $y=0$ to $y=1$
\item $x=\sqrt{y}$ from $y=0$ to $y=1$
\item $x=\frac{2}{3}(y^{2}+1)^{3/2}$ from $y=1$ to $y=3$
\item \techrequired $x=\tan y$ from $y=0$ to $y=\frac{3}{4}$
\item \techrequired $x=\cos^{2}y$ from $y=-\frac{\pi}{2}$ to $y=\frac{\pi}{2}$
\item \techrequired $x=4^{y}$ from $y=0$ to $y=2$
\item \techrequired $x=\ln(y)$ on $y=\frac{1}{e}$ to $y=e$
\end{enumerate}

For the following exercises, find the surface area of the volume generated when the following curves revolve around the $x\text{-axis}$. If you cannot evaluate the integral exactly, use your calculator to approximate it.

\begin{enumerate}[resume]
\item $y=\sqrt{x}$ from $x=2$ to $x=6$
\item $y=x^{3}$ from $x=0$ to $x=1$
\item $y=7x$ from $x=-1$ to $x=1$
\item \techrequired $y=\frac{1}{x^{2}}$ from $x=1$ to $x=3$
\item $y=\sqrt{4-x^{2}}$ from $x=0$ to $x=2$
\item $y=\sqrt{4-x^{2}}$ from $x=-1$ to $x=1$
\item $y=5x$ from $x=1$ to $x=5$
\item \techrequired $y=\tan x$ from $x=-\frac{\pi}{4}$ to $x=\frac{\pi}{4}$
\end{enumerate}

For the following exercises, find the surface area of the volume generated when the following curves revolve around the $y\text{-axis}$. If you cannot evaluate the integral exactly, use your calculator to approximate it.

\begin{enumerate}[resume]
\item $y=x^{2}$ from $x=0$ to $x=2$
\item $y=\frac{1}{2}x^{2}+\frac{1}{2}$ from $x=0$ to $x=1$
\item $y=x+1$ from $x=0$ to $x=3$
\item \techrequired $y=\frac{1}{x}$ from $x=\frac{1}{2}$ to $x=1$
\item $y=\sqrt[3]{x}$ from $x=1$ to $x=27$
\item \techrequired $y=3x^{4}$ from $x=0$ to $x=1$
\item \techrequired $y=\frac{1}{\sqrt{x}}$ from $x=1$ to $x=3$
\item \techrequired $y=\cos x$ from $x=0$ to $x=\frac{\pi}{2}$
\item The base of a lamp is constructed by revolving a quarter circle $y=\sqrt{2x-x^{2}}$ around the $y\text{-axis}$ from $x=1$ to $x=2$, as seen here. Create an integral for the surface area of this curve and compute it. \\ \textit{[Figure: This figure is a surface. It is half of a torus created by rotating the curve y=squareroot(2x-x\textasciicircum{}2) about the x-axis.]}
\item A light bulb is a sphere with radius $1/2$ in. with the bottom sliced off to fit exactly onto a cylinder of radius $1/4$ in. and length $1/3$ in., as seen here. The sphere is cut off at the bottom to fit exactly onto the cylinder, so the radius of the cut is $1/4$ in. Find the surface area (not including the top or bottom of the cylinder). \\ \textit{[Figure: This figure has two images. The first is a sphere on top of a cylinder. The second is a lightbulb.]}
\item \techrequired A lampshade is constructed by rotating $y=1/x$ around the $x\text{-axis}$ from $y=1$ to $y=2$, as seen here. Determine how much material you would need to construct this lampshade---that is, the surface area---accurate to four decimal places. \\ \textit{[Figure: This figure has two images. The first is similar to a frustum of a cone with edges bending inwards. The second is a lamp shade.]}
\item \techrequired An anchor drags behind a boat according to the function $y=24e^{-x/2}-24$, where $y$ represents the depth beneath the boat and $x$ is the horizontal distance of the anchor from the back of the boat. If the anchor is $23$ ft below the boat, how much rope do you have to pull to reach the anchor? Round your answer to three decimal places.
\item \techrequired You are building a bridge that will span $10$ ft. You intend to add decorative rope in the shape of $y=5|\sin((x\pi)/5)|$, where $x$ is the distance in feet from one end of the bridge. Find out how much rope you need to buy, measured in a whole number of feet.
\end{enumerate}

For the following exercises, find the exact arc length for the following problems over the given interval.

\begin{enumerate}[resume]
\item $y=\ln(\sin x)$ from $x=\pi/4$ to $x=(3\pi)/4$. (Hint: Recall trigonometric identities.)
\item \techrequired Draw graphs of $y=x^{2}$, $y=x^{6}$, and $y=x^{10}$. For $y=x^{n}$, as $n$ increases, formulate a prediction on the arc length from $(0,0)$ to $(1,1)$. Now, compute the lengths of these three functions and determine whether your prediction is correct.
\item Compare the lengths of the parabola $x=y^{2}$ and the line $x=by$ from $(0,0)$ to $(b^{2},b)$ as $b$ increases. What do you notice?
\item Solve for the length of $x=y^{2}$ from $(0,0)$ to $(1,1)$. Show that $x=(1/2)y^{2}$ from $(0,0)$ to $(2,2)$ is twice as long. Graph both functions and explain why this is so.
\item \techrequired Which is longer between $(1,1)$ and $(2,1/2)$: the hyperbola $y=1/x$ or the graph of $x+2y=3$?
\item Explain why the surface area is infinite when $y=1/x$ is rotated around the $x\text{-axis}$ for $1\le x<\infty$, but the volume is finite.
\end{enumerate}

\end{document}
